\SourceLeaf{163}
\SixthArticle{31}
Functio $\Psi z$ aeque fere memorabilis est atque functio $\Pi z$, quapropter insigniores relationes ad illam spectantes hic colligemus. E differentiatione aequationis 44 fit
\[
\Psi(z+1)=\Psi z+\frac{1}{z+1}\tag*{[63]}
\]
unde
\[
\Psi(z+n)=\Psi z+\frac{1}{z+1}+\frac{1}{z+2}+\frac{1}{z+3}+\text{etc.}+\frac{1}{z+n}\tag*{[64]}
\]
Huius adiumento a valoribus minoribus ipsius $z$ ad maiores progredi, vel a maioribus ad minores regredi licet: pro valoribus maioribus positivis ipsius $z$ functionis valores numerici satis commode per formulas sequentes e differentiatione aequationum 58, 59 oriundas computantur, de quibus tamen eadem sunt tenenda, quae in art. 29 circa formulas 58 et 59 monuimus:
\[
\Psi z=\log z+\frac{1}{2z}-\frac{\mathfrak{A}}{2zz}+\frac{\mathfrak{B}}{4z^4}-\frac{\mathfrak{C}}{6z^6}+\text{etc.}\tag*{[65]}
\]
\[
\Psi\left(z-\tfrac12\right)=\log z+\frac{\mathfrak{A}}{2\cdot2zz}-\frac{7\mathfrak{B}}{4\cdot8z^4}+\frac{31\mathfrak{C}}{6\cdot32z^6}-\text{etc.}\tag*{[66]}
\]
Ita pro $z=10$ computavimus
\[
\Psi z=2{,}3517525890\ 6672110764\ 743
\]
unde regredimur ad
\[
\Psi0=-0{,}5772156649\ 0153286060\ 653\HistoricalNote{18}
\]
Pro valore integro positivo ipsius $z$ fit generaliter
\[
\Psi z=\Psi0+1+\tfrac12+\tfrac13+\text{etc.}+\frac1z\tag*{[67]}
\]
Pro valore integro negativo autem manifesto $\Psi z$ fit quantitas infinite magna.
\SourceLeaf{164}
\SixthArticle{32}
Formula 55 nobis suppeditat $\log\Pi(-z)+\log\Pi(z-1)=\log\pi-\log\sin z\pi$, unde fit per differentiationem
\[
\Psi(-z)-\Psi(z-1)=\pi\operatorname{cotang}z\pi\tag*{[68]}
\]
Et quum e definitione functionis $\Psi$ generaliter habeatur
\[
\Psi x-\Psi y=-\frac1{x+1}+\frac1{y+1}-\frac1{x+2}+\frac1{y+2}-\frac1{x+3}+\text{etc.}\tag*{[69]}
\]
oritur series nota
\[
\pi\operatorname{cotang}z\pi=\frac1z-\frac1{1-z}+\frac1{1+z}-\frac1{2-z}+\frac1{2+z}-\frac1{3-z}+\text{etc.}
\]
Simili modo e differentiatione formulae 57 prodit
\[
\begin{aligned}
&\Psi z+\Psi\left(z-\frac1n\right)+\Psi\left(z-\frac2n\right)+\text{etc.}+\Psi\left(z-\frac{n-1}{n}\right)\\[4pt]
&=n\Psi nz-n\log n
\end{aligned}
\tag*{[70]}\]
adeoque statuendo $z=0$
\[
\begin{aligned}
&\Psi\left(-\frac1n\right)+\Psi\left(-\frac2n\right)+\Psi\left(-\frac3n\right)+\text{etc.}+\Psi\left(-\frac{n-1}{n}\right)\\[4pt]
&=(n-1)\Psi0-n\log n
\end{aligned}
\tag*{[71]}\]
Ita e. g. habetur
\[
\begin{aligned}
&\Psi\left(-\tfrac12\right)=\Psi0-2\log2=-1{,}9635100260\ 2142347944\ 099,\\[4pt]
&\quad\text{unde porro}
\end{aligned}
\]
\[
\Psi\tfrac12=+0{,}0364899739\ 7857652055\ 901.
\]
\SixthArticle{33}
Sicuti in art. praec. $\Psi\left(-\tfrac12\right)$ ad $\Psi0$ et logarithmum reduximus, ita generaliter $\Psi\left(-\frac mn\right)$, designantibus $m,n$ integros, quorum minor $m$, ad $\Psi0$ et logarithmos reducemus. Statuamus $\frac{2\pi}{n}=\omega$, sitque $\varphi$ alicui angulorum $\omega,2\omega,3\omega\ .\ .\ .\ .\ (n-1)\omega$ aequalis; unde $1=\cos n\varphi=\cos2n\varphi=\cos3n\varphi$ etc., $\cos\varphi=\cos(n+1)\varphi=\cos(n+2)\varphi$ etc., $\cos2\varphi=\cos(n+2)\varphi$ etc., nec non $\cos\varphi+\cos2\varphi+\cos3\varphi+\text{etc.}+\cos(n-1)\varphi+1=0$. Habemus itaque
\[
\begin{aligned}
&\cos\varphi\cdot\Psi\frac{1-n}{n}=-n\cos\varphi+\cos\varphi\cdot\log2\\[4pt]
&-\frac{n}{n+1}\cos(n+1)\varphi+\cos\varphi\cdot\log\tfrac32-\text{etc.}
\end{aligned}
\]
\[
\begin{aligned}
&\cos2\varphi\cdot\Psi\frac{2-n}{n}=-\frac n2\cos2\varphi+\cos2\varphi\cdot\log2\\[4pt]
&-\frac{n}{n+2}\cos(n+2)\varphi+\cos2\varphi\cdot\log\tfrac32-\text{etc.}
\end{aligned}
\]
\[
\begin{aligned}
&\cos3\varphi\cdot\Psi\frac{3-n}{n}=-\frac n3\cos3\varphi+\cos3\varphi\cdot\log2\\[4pt]
&-\frac{n}{n+3}\cos(n+3)\varphi+\cos3\varphi\cdot\log\tfrac32-\text{etc.}
\end{aligned}
\]
etc. usque ad
\SourceLeaf{165}
\[
\begin{aligned}
&\cos(n-1)\varphi\cdot\Psi\left(-\frac1n\right)
\\[4pt]
&=-\frac{n}{n-1}\cos(n-1)\varphi+\cos(n-1)\varphi\cdot\log2\\[4pt]
&-\frac{n}{2n-1}\cos(2n-1)\varphi
+\cos(n-1)\varphi\cdot\log\tfrac32-\text{etc.}
\end{aligned}
\]
\[
\Psi0=-\frac nn\cos n\varphi+\log2-\frac{n}{2n}\cos2n\varphi+\log\tfrac32-\text{etc.}
\]
atque per \textit{summationem}
\[
\begin{aligned}
&\cos\varphi\cdot\Psi\frac{1-n}{n}+\cos2\varphi\cdot\Psi\frac{2-n}{n}\\[4pt]
&+\cos3\varphi\cdot\Psi\frac{3-n}{n}+\text{etc.}\\[4pt]
&+\cos(n-1)\varphi\cdot\Psi\left(-\frac1n\right)+\Psi0
\\[4pt]
&=-n(\cos\varphi+\tfrac12\cos2\varphi+\tfrac13\cos3\varphi+\tfrac14\cos4\varphi+\text{etc. in infin.})
\end{aligned}
\]
Sed habetur generaliter, pro valore ipsius $x$ unitate non maiori,
\[
\begin{aligned}
&\log(1-2x\cos\varphi+xx)\\[4pt]
&=-2(x\cos\varphi+\tfrac12xx\cos2\varphi+\tfrac13x^3\cos3\varphi+\text{etc.})
\end{aligned}
\]
quae quidem series facile sequitur ex evolutione $\log(1-rx)+\log\left(1-\frac xr\right)$, denotante $r$ quantitatem $\cos\varphi+\BareRoot{-1}\cdot\sin\varphi$. Hinc fit aequatio praecedens
\[
\begin{aligned}
&\cos\varphi\cdot\Psi\frac{1-n}{n}+\cos2\varphi\cdot\Psi\frac{2-n}{n}\\[4pt]
&+\cos3\varphi\cdot\Psi\frac{3-n}{n}+\text{etc.}\\[4pt]
&+\cos(n-1)\varphi\cdot\Psi\left(-\frac1n\right)
\\[4pt]
&=-\Psi0+\tfrac12n\log(2-2\cos\varphi)
\end{aligned}
\tag*{[72]}\]
Statuatur in hac aequatione deinceps $\varphi=\omega$, $\varphi=2\omega$, $\varphi=3\omega$ etc. usque ad $\varphi=(n-1)\omega$, multiplicentur singulae hae aequationes ordine suo per $\cos m\omega$, $\cos2m\omega$, $\cos3m\omega$ etc. usque ad $\cos(n-1)m\omega$, productorumque aggregato adiiciatur aequatio 71
\[
\begin{aligned}
&\Psi\frac{1-n}{n}+\Psi\frac{2-n}{n}+\Psi\frac{3-n}{n}+\text{etc.}+\Psi\left(-\frac1n\right)\\[4pt]
&=(n-1)\Psi0-n\log n
\end{aligned}
\]
Quodsi iam perpenditur, esse
\[
\begin{aligned}
&1+\cos m\omega\cdot\cos k\omega+\cos2m\omega\cdot\cos2k\omega\\[4pt]
&+\cos3m\omega\cdot\cos3k\omega
+\text{etc.}+\cos(n-1)m\omega\cdot\cos(n-1)k\omega=0
\end{aligned}
\]
denotante $k$ aliquem numerorum $1,2,3\ .\ .\ .\ .\ (n-1)$ exceptis his duobus $m$ atque $n-m$, pro quibus summa illa fit $=\tfrac12n$, patebit, ex summatione illarum aequationum prodire, post divisionem per $\frac n2$,
\[
\begin{aligned}
&\Psi\left(-\frac mn\right)+\Psi\left(-\frac{n-m}{m}\right)=
\\[4pt]
&2\Psi0-2\log n+\cos m\omega\cdot\log(2-2\cos\omega)\\[4pt]
&+\cos2m\omega\cdot\log(2-2\cos2\omega)
\\[4pt]
&+\cos3m\omega\cdot\log(2-2\cos3\omega)+\text{etc.}\\[4pt]
&+\cos(n-1)m\omega\cdot\log(2-2\cos(n-1)\omega)
\end{aligned}
\tag*{[73]}\]
Manifesto terminus ultimus huius aequationis fit $=\cos m\omega\cdot\log(2-2\cos\omega)$, pen\SourceLeaf{166}ultimus $=\cos 2m\omega.\log(2-2\cos 2\omega)$ etc., ita ut bini termini semper sint aequales, excepto, si $n$ est par, termino singulari $\cos\frac{n}{2}.m\omega\log(2-2\cos\frac{n}{2}\omega)$, qui fit $=+2\log 2$ pro $m$ pari, vel $=-2\log 2$ pro $m$ impari. Combinando iam cum aequatione 73 hanc
\[
\Psi\left(-\frac{m}{n}\right)-\Psi\left(-\frac{n-m}{n}\right)=\pi\operatorname{cotang}\frac{m}{n}\pi
\]
habemus, pro valore impari ipsius $n$, siquidem $m$ est integer positivus minor quam $n$
\[
\begin{aligned}
&{}\text{[74]}\qquad \Psi\left(-\frac{m}{n}\right)
\\[4pt]
&=\Psi 0+\tfrac12\pi\operatorname{cotang}\frac{m\pi}{n}-\log n\\[4pt]
&+\cos\frac{2m\pi}{n}.\log\left(2-2\cos\frac{2\pi}{n}\right)
\\[4pt]
&+\cos\frac{4m\pi}{n}.\log\left(2-2\cos\frac{4\pi}{n}\right)\\[4pt]
&+\cos\frac{6m\pi}{n}.\log\left(2-2\cos\frac{6\pi}{n}\right)+\text{etc.}
\\[4pt]
&+\cos\frac{(n-1)m\pi}{n}.\log\left(2-2\cos\frac{(n-1)\pi}{n}\right)
\end{aligned}
\]
Pro valore pari ipsius $n$ autem
\[
\begin{aligned}
&{}\text{[75]}\qquad \Psi\left(-\frac{m}{n}\right)
\\[4pt]
&=\Psi 0+\tfrac12\pi\operatorname{cotang}\frac{m\pi}{n}-\log n\\[4pt]
&+\cos\frac{2m\pi}{n}\log\left(2-2\cos\frac{2\pi}{n}\right)
\\[4pt]
&+\cos\frac{4m\pi}{n}\log\left(2-2\cos\frac{4\pi}{n}\right)+\text{etc.}\\[4pt]
&+\cos\frac{(n-2)m\pi}{n}\log\left(2-2\cos\frac{(n-2)\pi}{n}\right)
\pm\log 2
\end{aligned}
\]
ubi signum superius valet pro $m$ pari, inferius pro impari. Ita e.~g. invenitur
\[
\begin{aligned}
&\Psi\left(-\tfrac14\right)=\Psi 0+\tfrac12\pi-3\log 2, 
\\[4pt]
&\Psi\left(-\tfrac34\right)=\Psi 0-\tfrac12\pi-3\log 2
\\[4pt]
&\Psi\left(-\tfrac13\right)=\Psi 0+\tfrac12\pi\BareRoot{\tfrac13}-\tfrac32\log 3, 
\\[4pt]
&\Psi\left(-\tfrac23\right)=\Psi 0-\tfrac12\pi\BareRoot{\tfrac13}-\tfrac22\log 3
\end{aligned}
\]
Ceterum combinatis his aequationibus cum aequatione 64 sponte patet, $\Psi z$ generaliter pro \emph{quovis valore rationali} ipsius $z$, positivo seu negativo per $\Psi 0$ atque logarithmos determinari posse, quod theorema sane maxime est memorabile.

\SixthArticle{34}
Quum, per art.~28, $\Pi\lambda$ sit valor integralis $\int y^\lambda e^{-y}\mathrm{d}y$, ab $y=0$ usque ad $y=\infty$, siquidem $\lambda+1$ est quantitas positiva, fit differentiando secundum $\lambda$
\[
\frac{\mathrm{d}\Pi\lambda}{\mathrm{d}\lambda}
=\frac{\mathrm{d}\int y^\lambda e^{-y}\mathrm{d}y}{\mathrm{d}\lambda}
=\int y^\lambda e^{-y}\log y\,\mathrm{d}y
\]
sive
\[
\begin{aligned}
&[76]\qquad \Pi\lambda.\Psi\lambda=\int y^\lambda e^{-y}\log y.\mathrm{d}y,\\[4pt]
&\quad\text{ab }y=0\text{ usque ad }y=\infty
\end{aligned}
\]
\par\SourceLeaf{167}
{\raggedright\parindent=1.2em Generalius statuendo $y=z^\alpha$, $\alpha\lambda+\alpha-1=\beta$, valor integralis $\int z^\beta e^{-z^\alpha}\log z\,\mathrm{d}z$, a $z=0$ usque ad $z=\infty$, fit\par}
\[
\begin{aligned}
&=\frac{1}{\alpha\alpha}\Pi\left(\frac{\beta+1}{\alpha}-1\right).\Psi\left(\frac{\beta+1}{\alpha}-1\right)
\\[4pt]
&=\frac{1}{\alpha(\beta+1)}\Pi\frac{\beta+1}{\alpha}.\Psi\frac{\beta+1}{\alpha}
-\frac{1}{(\beta+1)^2}\Pi\frac{\beta+1}{\alpha}
\end{aligned}
\]
siquidem simul $\beta+1$ atque $\alpha$ sunt quantitates positivae, vel aequalis eidem quantitati cum signo opposito, si utraque $\beta+1$, $\alpha$ est negativa.

\SixthArticle{35}
At non solum productum $\Pi\lambda.\Psi\lambda$, verum etiam ipsa functio $\Psi\lambda$ per integrale determinatum exhiberi potest. Designante $k$ integrum positivum, patet valorem integralis $\int\frac{x^\lambda-x^{\lambda+k}}{1-x}.\mathrm{d}x$, ab $x=0$ usque ad $x=1$ esse
\[
=\frac{1}{\lambda+1}+\frac{1}{\lambda+2}+\frac{1}{\lambda+3}+\text{etc.}+\frac{1}{\lambda+k}
\]
Porro quum valor integralis $\int\left(\frac{1}{1-x}-\frac{kx^{k-1}}{1-x^k}\right)\mathrm{d}x$ generaliter sit $=\mathrm{Const.}+\log\frac{1-x^k}{1-x}$, idem inter limites $x=0$ atque $x=1$ erit $=\log k$, unde patet, valorem integralis $S=\int\left(\frac{1-x^\lambda+x^{\lambda+k}}{1-x}-\frac{kx^{k-1}}{1-x^k}\right)\mathrm{d}x$ inter eosdem limites esse
\[
=\log k-\frac{1}{\lambda+1}-\frac{1}{\lambda+2}-\frac{1}{\lambda+3}-\text{etc.}-\frac{1}{\lambda+k}
\]
quam expressionem denotabimus per $\Omega$. Discerpamus integrale $S$ in duas partes
\[
\int\left(\frac{1-x^\lambda}{1-x}\right)\mathrm{d}x+\int\left(\frac{x^{\lambda+k}}{1-x}-\frac{kx^{k-1}}{1-x^k}\right)\mathrm{d}x
\]
Pars prima $\int\frac{1-x^\lambda}{1-x}.\mathrm{d}x$, statuendo $x=y^k$ mutatur in
\[
\int\frac{ky^{k-1}-ky^{\lambda k+k-1}}{1-y^k}\mathrm{d}y
\]
unde sponte patet, illius valorem ab $x=0$ usque ad $x=1$, aequalem esse valori integralis
\[
\int\frac{kx^{k-1}-kx^{\lambda k+k-1}}{1-x^k}\mathrm{d}x
\]
inter eosdem limites, quum manifesto literam $y$ sub hac restrictione in $x$ mutare liceat. Hinc fit integrale $S$, inter eosdem limites
\[
=\int\left(\frac{x^{\lambda+k}}{1-x}-\frac{kx^{\lambda k+k-1}}{1-x^k}\right)\mathrm{d}x
\]
Hoc vero integrale, statuendo $x^k=z$, transit in
\SourceLeaf{168}
\[
\int\left(\frac{z^{\frac{\lambda+1}{k}}}{k(1-z)^{\frac{1}{k}}}-\frac{z^\lambda}{1-z}\right)\mathrm{d}z
\]
quod itaque inter limites $z=0$ atque $z=1$ sumtum aequale est ipsi $\Omega$. Sed crescente $k$ in infinitum, limes ipsius $\Omega$ est $\Psi\lambda$, limes ipsius $\frac{\lambda+1}{k}$ est $0$. limes ipsius $k(1-z)^{\frac{1}{k}}$ vero est $\log\frac{1}{z}$ sive $-\log z$. Quare habemus
\[
\begin{aligned}
&[77]\qquad \Psi\lambda\\[4pt]
&=\int\left(\frac{1}{\log\frac{1}{z}}-\frac{z^\lambda}{1-z}\right)\mathrm{d}z
=\int\left(-\frac{1}{\log z}-\frac{z^\lambda}{1-z}\right)\mathrm{d}z
\end{aligned}
\]
a $z=0$ usque ad $z=1$.

\SixthArticle{36}
Integralia determinata, per quae supra expressae sunt functiones $\Pi\lambda$, $\Pi\lambda.\Psi\lambda$, restringere oportuit ad valores ipsius $\lambda$ tales, ut $\lambda+1$ evadat quantitas positiva: haec restrictio ex ipsa deductione demanavit, reveraque facile perspicitur, pro aliis valoribus ipsius $\lambda$ illa integralia semper fieri infinita, etiamsi functiones $\Pi\lambda$, $\Pi\lambda.\Psi\lambda$ finitae manere possint. Veritati formula 77 certo eadem conditio subesse debet, ut $\lambda+1$ sit quantitas positiva (alioquin enim integrale certo infinitum evadit, etiamsi functio $\Psi\lambda$ maneat finita): sed deductio formulae primo aspectu generalis nullique restrictioni obnoxia esse videtur. Sed propius attendenti facile patebit, ipsi analysi, per quam formula eruta est, hanc restrictionem iam inesse. Scilicet tacite supposuimus, integrale $\int\frac{1-x^\lambda}{1-x}\mathrm{d}x$ cui aequale $\int\frac{kx^{k-1}-kx^{\lambda k+k-1}}{1-x^k}\mathrm{d}x$ substituimus, habere valorem \emph{finitum}, quae conditio requirit, ut $\lambda+1$ sit quantitas positiva. Ex analysi nostra quidem sequitur, haec duo integralia semper esse aequalia, si hoc extendatur ab $x=0$ usque ad $x=1-\omega$, illud ab $x=0$ usque ad $x=(1-\omega)^k$, quantumvis parva sit quantitas $\omega$, modo non sit $=0$: sed hoc non obstante in casu eo, ubi $\lambda+1$ non est quantitas positiva, duo integralia ab $x=0$ usque ad \emph{eundem} terminum $x=1-\omega$ extensa neutiquam ad aequalitatem convergunt, sed potius tunc ipsorum differentia, decrescente $\omega$ in infinitum, in infinitum crescet. Hocce exemplum monstrat, quanta circumspicientia opus sit in tractandis quantitatibus infinitis, quae in ratiociniis analyticis nostro iudicio eatenus tantum sunt admittendae, quatenus ad theoriam limitum reduci possunt.
\SourceLeaf{169}
\SixthArticle{37}
Statuendo in formula 77, $z=e^{-u}$, patet, illam etiam ita exhiberi posse
\[
\begin{aligned}
&\Psi\lambda=-\int\left(\frac{e^{-u}}{u}-\frac{e^{-u\lambda-u}}{1-e^{-u}}\right)\mathrm{d}u,
\\[4pt]
&\quad\text{ab }u=\infty\text{ usque ad }u=0,\text{ i. e.}
\end{aligned}
\]
\[
\begin{aligned}
&\text{[78]}\qquad
\Psi\lambda=\int\left(\frac{e^{-u}}{u}-\frac{e^{-\lambda u}}{e^u-1}\right)\mathrm{d}u,
\\[4pt]
&\quad\text{ab }u=0\text{ usque ad }u=\infty.
\end{aligned}
\]
(Perinde valor ipsius $\Pi\lambda$ in art. 28 allatus, mutatur statuendo
$e^{-y}=\upsilon$, in sequentem
\[
\Pi\lambda=\int\left(\log\frac{1}{\upsilon}\right)^\lambda\mathrm{d}\upsilon,
\quad\text{a }\upsilon=0\text{ usque ad }\upsilon=1)
\]
Porro patet e formula 77, esse
\[
\text{[79]}\qquad
\Psi\lambda-\Psi\mu=\int\frac{z^\mu-z^\lambda}{1-z}\mathrm{d}z,
\quad\text{a }z=0\text{ usque ad }z=1
\]
ubi praeter $\lambda+1$ etiam $\mu+1$ debet esse quantitas positiva.

Statuendo in eadem formula 77, $z=u^\alpha$, designante $\alpha$ quantitatem
positivam, fit
\[
\begin{aligned}
&\Psi\lambda=\int\left(-\frac{u^{\alpha-1}}{\log u}
-\frac{\alpha u^{\alpha\lambda+\alpha-1}}{1-u^\alpha}\right)\mathrm{d}u,
\\[4pt]
&\quad\text{ab }u=0\text{ usque ad }u=1
\end{aligned}
\]
et quum perinde statui possit, pro valore positivo ipsius $\beta$,
\[
\Psi\lambda=\int\left(-\frac{u^{\beta-1}}{\log u}
-\frac{\beta u^{\beta\lambda+\beta-1}}{1-u^\beta}\right)\mathrm{d}u
\]
patet, fieri
\[
0=\int\left(\frac{u^{\alpha-1}-u^{\beta-1}}{\log u}
+\frac{\alpha u^{\alpha\lambda+\alpha-1}}{1-u^\alpha}
-\frac{\beta u^{\beta\lambda+\beta-1}}{1-u^\beta}\right)\mathrm{d}u
\]
sive
\[
\int\frac{u^{\alpha-1}-u^{\beta-1}}{\log u}\mathrm{d}u
=\int\left(\frac{\beta u^{\beta\lambda+\beta-1}}{1-u^\beta}
-\frac{\alpha u^{\alpha\lambda+\alpha-1}}{1-u^\alpha}\right)\mathrm{d}u
\]
integralibus semper ab $u=0$ usque ad $u=1$ extensis. Sed ponendo $\lambda=0$,
integrale posterius \emph{indefinite} assignari potest; est scilicet
$=\log\frac{1-u^\alpha}{1-u^\beta}$, si evanescere debet pro $u=0$;
quare quum pro $u=1$ statuere oporteat
$\frac{1-u^\alpha}{1-u^\beta}=\frac{\alpha}{\beta}$,
erit integrale
$\log\frac{\alpha}{\beta}=\int\frac{u^{\alpha-1}-u^{\beta-1}}{\log u}\mathrm{d}u$,
ab $u=0$ usque ad $u=1$, quod theorema olim ab ill. \textsc{Euler} per alias
methodos erutum est.
\par\medskip\noindent\hfill\rule{.25\linewidth}{.4pt}\hfill\par
